\documentclass[12pt,a4paper]{article}

\usepackage[T1]{fontenc}
\usepackage[utf8]{inputenc}
\usepackage[brazil]{babel}
\usepackage{lmodern}
\usepackage{geometry}
\usepackage{hyperref}
\usepackage{enumitem}
\usepackage{array}

\geometry{margin=2.5cm}
\hypersetup{
  colorlinks=true,
  linkcolor=blue,
  urlcolor=blue
}

\title{Plano de Aula: Procedimentos Armazenados}
\date{}

\begin{document}
\maketitle

\noindent
\textbf{Duração:} 2 horas (120 minutos)\\
\textbf{Público-alvo:} Iniciantes\\
\textbf{Recursos Necessários:} Laboratório com computadores, SGBD e cliente SQL, projetor (opcional), quadro e marcadores.

\section{Objetivos da Aula}
\begin{itemize}[leftmargin=*,itemsep=0.25em]
  \item Compreender o que são procedimentos armazenados e quando utilizá-los.
  \item Criar procedimentos com parâmetros de entrada e saída (quando aplicável).
  \item Encapsular lógica de negócio simples no banco com validações e consultas.
  \item Relacionar procedures com reutilização, segurança e desempenho.
\end{itemize}

\section{Metodologia}
Exposição curta com analogias e prática em laboratório. Os alunos criam procedures pequenas para o domínio da clínica (ex.: cadastrar pet, listar pets de um cliente) e testam com casos de entrada diferentes.

\section{Cronograma e Conteúdo Detalhado (120 Minutos)}
\subsection*{Parte 1: Conceito e Motivação (20 minutos)}
\begin{itemize}[leftmargin=*,itemsep=0.35em]
  \item O que é uma procedure: ``uma receita pronta que o SGBD executa''.
  \item Analogia: ``macro'' que recebe parâmetros e executa passos em sequência.
  \item Diferença entre executar SQL direto e encapsular em procedure.
\end{itemize}

\subsection*{Parte 2: Sintaxe e Componentes (25 minutos)}
\begin{itemize}[leftmargin=*,itemsep=0.35em]
  \item CREATE PROCEDURE (ou equivalente do SGBD adotado).
  \item Parâmetros (IN/OUT/INOUT, quando aplicável).
  \item Variáveis, controle simples (IF) e retorno de resultados.
\end{itemize}

\subsection*{Parte 3: Exemplo Guiado (25 minutos)}
\begin{itemize}[leftmargin=*,itemsep=0.35em]
  \item Criar procedure \textbf{CadastrarPet} (inserir pet para um cliente existente).
  \item Validar FK: se o cliente não existir, sinalizar erro (quando suportado) ou retornar mensagem.
  \item Criar procedure \textbf{ListarPetsDoCliente} (consulta com JOIN).
\end{itemize}

\subsection*{Parte 4: Atividade em Duplas --- Biblioteca de Procedures (40 minutos)}
\begin{itemize}[leftmargin=*,itemsep=0.35em]
  \item Cada dupla cria 2 procedures:
  \begin{itemize}[leftmargin=*,itemsep=0.25em]
    \item Uma de \textbf{escrita} (INSERT/UPDATE) com validação.
    \item Uma de \textbf{leitura} (SELECT) com filtro por parâmetro.
  \end{itemize}
  \item Testar com pelo menos 3 casos (entrada válida, inválida e limite).
  \item Entrega: scripts e evidências (resultados no console).
\end{itemize}

\subsection*{Parte 5: Encerramento (10 minutos)}
\begin{itemize}[leftmargin=*,itemsep=0.35em]
  \item Revisar: vantagens, riscos e boas práticas.
  \item Ponte: ``Próxima aula: regras automáticas --- asserções e gatilhos.''
\end{itemize}

\end{document}

